% page 530 of the facsimile (image p-529 of the scan)
% block 160.0 x 254.9 mm ; left margin 8.0 mm
\pgc{-12.700mm}{0.000mm}{
\l{\cel{3}522}
\l{}
\l{\sou{sexa-}\cel{1}.Prefixo ciencala : \textquotedbl{}qua havas sis...\textquotedbl{}}
\l{}
\l{\sou{sexagesimo}. (liturgio katol.) La dio sisadekesma ante la}
\l{\cel{3}oktavo di pasko. -}
\l{}
\l{\sou{sexto}. (muziko). Intervalo qua kontenas sis noti - de \textquotedbl{}do\textquotedbl{}}
\l{\cel{3}a \textquotedbl{}la\textquotedbl{}. - II. (\sou{skermo}) Pareo per la espado vers-supre, en}
\l{\cel{3}la direciono adextere, kun la karpo vers-ciele. - DEFIRS}
\l{}
\l{\sou{sextanto}. I. (\sou{geom.)}\cel{2}La parto sisima di cirklokurvo (arko}
\l{\cel{3}de 60 gradi). - II. Instrumento di qua la parto precipua}
\l{\cel{3}esas limbo gradizita, longa ye la sisimo di cirklokurvo}
\l{\cel{3}t.e. 60 gradi, e qua uzesas por mezurar la anguli til 60}
\l{\cel{3}gradi por determinar la disti di astri, la situeso di}
\l{\cel{3}navo sur maro, e c. - DEFIRS.}
\l{}
\l{\sou{sexteto}. (\sou{muziko)}.\cel{2}Muzikajo por sis instrumenti o por sis}
\l{\cel{3}voci (kun o sen akompano). - DEFIRS.}
\l{}
\l{\sou{sextiliono}. (\sou{matem.)}\cel{2}1036}
\l{}
\l{\sou{sexuo}. I. Karaktero organala naturala, qua interdistingas}
\l{\cel{3}la maskulo e la femino. - II. Ensemblo de la karakteri}
\l{\cel{3}organala, pensala, morala, qui interdistingas la homulo}
\l{\cel{3}e la homino. -\cel{1}\sou{III}. La ensemblo de la individui di ta}
\l{\cel{3}o ca sexuo. - DEFIRS.}
\l{}
\l{\sou{sezamo}. (\sou{bot.}) Planto oleifera ek la regioni estala. - L.}
\l{\cel{3}\sou{sesamum}.}
\l{}
\l{\sou{sezono}. Singla ek la quar dividuri inter-egala di la yaro,}
\l{\cel{3}determinita segun la quanto de tempo quan spensas Suno}
\l{\cel{3}por irar de solstico ad equinoxo, e de equinoxo a sols-}
\l{\cel{3}tico, e segun la varii temperaturala qui korespondas a}
\l{\cel{3}li. - DEFR.}
\l{}
\l{\sou{sfacelo.}\cel{2}(\sou{medic}.) Gangreno qua atakas la tota parto dika}
\l{\cel{3}di membro o di organo plura-tisua. -}
\l{}
\l{\sou{sfenoida.(anat.}) Osto, ye la bazo di la kranio, qua kontri-}
\l{\cel{3}butas a la formaco di la orbiti, di la nazo-kavaji.}
\l{}
\l{\sou{sfero}. (\sou{geom.})\cel{2}Solido qua havas kom limito surfaco kurva}
\l{\cel{3}di qua omna punti distas egale de\cel{2}punto interna, la}
\l{\cel{3}\textquotedbl{}centro\textquotedbl{}. - DEFIS.}
\l{}
\l{\sou{sferoido}. (\sou{geom.}) Solido di qua la formo esas preske sfera.}
\l{\cel{3}- DEFIS.}
\l{}
\l{\sou{sfintero}. (\sou{anat.)}\cel{2}Muskulo cirklokurva qua klozas la orifico}
\l{\cel{3}di ula kavaji naturala di la homo-korpo (anuso, veziko).}
\l{\cel{3}- DEFIS.}
\l{}
\l{\sou{sfinxo}.I. (\sou{mitol}.) Monstro qua havas kapo e manui quale homo,}
\l{\cel{3}korpo quale leono, ali quale aglo, e qua devoris la per-}
\l{\cel{3}soni qui ne divinis la enigmati quin lu propozis. - II.}
\l{\cel{3}(\sou{metaf}.) Persono qua ne lasas sua pensi divinesar. - DEFIRS.}
}
